\section{Where the cost of index calculus lives}\label{sec:model}

An improvement is only useful if it moves a term that matters. This section
writes the cost of one index-calculus solve as a sum of terms, and maps each
lever of the programme to the term it acts on.

\subsection{The pipeline}

Fix a factor base $\Fb\subset E(\FF_q)$ of size $N$. A \emph{relation} is a
decomposition $aG+bQ=P_1+\cdots+P_m$ with every $P_i\in\pm\Fb$. Once about
$N$ independent relations are known, linear algebra over $\ZZ/r$ yields
$\log_G Q$. A heuristic count is enough to place the levers. A uniformly
random point decomposes into $m$ base points with probability about
$N^m/(m!\,r)$. Collecting $N$ relations therefore needs about
$m!\,r/N^{m-1}$ trials. The cold cost of one target is then
\begin{equation}\label{eq:cost}
  T \;\approx\;
  \underbrace{T_{\mathrm{setup}}(N)}_{\text{base, tables}}
  \;+\;
  \underbrace{\frac{m!\,r}{N^{m-1}}\;\frac{1}{\rho_{\mathrm{yield}}}\;
  C_{\mathrm{dec}}(N,m)}_{\text{relation collection}}
  \;+\;
  \underbrace{T_{\mathrm{LA}}(N,m)}_{\text{linear algebra}}
  \;+\;
  \underbrace{T_{\mathrm{target}}}_{\text{descent, check}}.
\end{equation}
Here $C_{\mathrm{dec}}$ is the cost of one decomposition attempt (the
\emph{oracle}). The factor $\rho_{\mathrm{yield}}\ge1$ collects everything
that makes a trial more productive than the heuristic: free relations,
large primes, symmetry, re-randomised targets. $T_{\mathrm{LA}}$ is roughly
$mN^2$ for sparse elimination. With rho's cost $\sqrt{\pi r/(2A)}$ on the
other side, the classical balance at $m=2$ with a constant-cost oracle gives
$N\approx r^{1/3}$ and $T\approx r^{2/3}$. That is why two-summand
prime-field index calculus loses ground to rho as $r^{1/6}$
(\S\ref{sec:distance}).

\subsection{The levers}

Table~\ref{tab:levers} lists, for each term of~\eqref{eq:cost}, the levers
the programme pulled and the sections that report them.

\begin{table}[t]
\centering\small
\caption[]{The terms of~\eqref{eq:cost} and the levers that act on them.}
\label{tab:levers}
\begin{tabularx}{\linewidth}{@{}l>{\raggedright}X>{\raggedright\arraybackslash}p{2.6cm}@{}}
\toprule
Term & Levers pulled & Section \\
\midrule
$C_{\mathrm{dec}}$ & root finding of $S_3$; pair and meet-in-the-middle
  tables; inherited F4, F5, F6-IC; gate-first order; harvest gates;
  $\FF_p$ F4; $D_4$ invariants & \S\ref{sec:oracle}, \S\ref{sec:families} \\
$\rho_{\mathrm{yield}}$, $N$ & compact Frobenius orbits; free relations;
  stride-renewed targets; greedy coverage; $n=83$ base selection; double
  large primes & \S\ref{sec:basis}, \S\ref{sec:families} \\
$T_{\mathrm{setup}}$ & stacked producer; cache-blocked pair table; u64 root
  batching & \S\ref{sec:setup} \\
$T_{\mathrm{LA}}$ & CADO-style filter and Wiedemann; structured Gaussian
  elimination; block Wiedemann; first-dependency stop & \S\ref{sec:setup},
  \S\ref{sec:families} \\
the exponent & genus-three transfer (cubic GHS); $k=4$ extension fields &
  \S\ref{sec:families} \\
\bottomrule
\end{tabularx}
\end{table}

\subsection{Which term dominates now}

Early in the programme, $C_{\mathrm{dec}}$ dominated. A Gr\"obner solve per
trial cost milliseconds to seconds. After root finding
(\S\ref{sec:oracle}), the oracle is a handful of field operations and table
probes, and the cost has moved elsewhere.
\begin{itemize}
\item On the Koblitz descent of the leaderboard's Table~A, reusable set-up
  is $98.5$--$99.7\%$ of cold $S$, and descent is $0.3$--$1.4\%$
  (\art{docs/ic/LEADERBOARD.md}, phase table).
\item On the $n=83$ cold pipeline with base K1107, a mean $113.755$\,s
  process splits into $47.166$\,s of set-up and $66.486$\,s of rank stage.
  The target stages take $0.099$\,s
  (\art{research/n83_cold_base_confirm_20261021/RESULT.md}).
\end{itemize}
So the relation-collection and set-up terms dominate. Their cost is set by
$N$ and by the yield per probe, not by the oracle. This is the main
structural fact behind the ledger. Oracle speed-ups of three to five orders
of magnitude moved the end-to-end gap by one order at most.
